\section{Setting Up to Write Together}

The advice of the last section assumes a team, and teams arrive with
unequal tooling comfort.
Some authors live in the terminal; others have never typed a git
command and never should.
Forcing one tool on everyone loses authors at both ends; hence we
offer a ladder, and each author picks a rung:
\begin{enumerate}
\item a reader, who watches the built PDF and mails comments; \item a
browser writer, who drafts in a shared Google Doc with help from
claude.ai; \item an Overleaf writer, who edits the LaTeX live in the
browser and watches the PDF re-render; \item a git writer, who clones
the repository, edits, and pushes; \item an agent runner, who works
with Claude Code (or similar) on their own machine, with the agent
reading and writing the shared source directly.
\end{enumerate}
Every rung is a full member of the team; the rungs differ only in
setup cost.

One piece of plumbing joins the ends of that ladder.
An Overleaf project is also a git remote, so the same paper has two
faces: a live browser editor for humans and an ordinary repository for
people and agents.
GitHub holds the source of truth; a push moves edits to Overleaf in
seconds; a pull collects whatever the browser writers typed.
Hence n people and n agents can write asynchronously, meeting only at
the merge.

Merges stay small because of two source-layout habits.
Firstly, every sentence ends its line, so git diffs and merges by the
sentence, and two edits collide only when they touch the same
sentence.
Secondly, prose lives in many small files (one per section) beneath a
thin root of include statements, so writers claim a section, not the
paper.
Review comments follow the same philosophy: a comment is a macro in
the source (who said it, what they said), listed by one make target
and resolved by deleting it.
Margin comments in web editors were rejected since they do not survive
the trip through git.

The last ritual is scheduling.
Before any shared writing session, one owner verifies that every
participant can already write to something: doc links delivered,
project invites accepted, push access proven by one trivial commit.
Access problems solved during a session cost n people an hour; solved
the day before, they cost one person a minute.
The prompts the team uses are versioned in the same repository as the
prose, so every rung (browser paste or agent invocation) runs the same
words.
(Aside: this paper was written with exactly this setup, one author
plus one agent, trading pushes through the bridge described above.)
